\section{Introduction}
\label{sec:intro}

As datacenter network links scale to 100--400\,Gb/s and RPC payloads shrink to tens or hundreds of bytes, the host--device interconnect has become a primary bottleneck for end-host networking~\cite{neugebauer2018pcie,agarwal2023hostcc,vuppalapati2024hostnet,sadok2023enso}. Standard descriptor-ring interfaces force every small outbound message through host memory: the CPU writes payloads and descriptors into host DRAM or LLC, rings a doorbell, and waits for the NIC to issue inbound PCIe DMA reads back across the root complex. That round-trip consumes host LLC capacity, steals memory controller (IMC) bandwidth from co-located compute, and adds microseconds of queuing delay~\cite{farshin2020ddio,liu2025ceio}.

A growing line of architectures and commercial devices therefore promotes \emph{MMIO direct push}---often realized as Programmed I/O (PIO), Mellanox/NVIDIA BlueFlame, Intel \texttt{MOVDIR64B} posted writes, or accelerator submission portals~\cite{kalia2016rdma,liew2025hotos,ruzhanskaia2024pio,inteldsa2023}. Rather than staging payloads in host RAM for the device to pull, a host core writes 64-byte cache lines directly into device-mapped MMIO space. In steady state, this mechanism delivers a striking isolation dividend: on our Intel Sapphire Rapids server, a single core pushing 175\,Mpps (11.2\,GB/s) of 64-byte posted writes directly into a Marvell OCTEON~10 CN106 SmartNIC's \texttt{BAR4} window degrades co-located 6-core STREAM bandwidth and L3/DRAM pointer-chase latency on the same socket by less than 1\%. Pushing the same message rate through a host-memory DMA ring degrades co-located STREAM bandwidth by 9.9\%--41.0\%.

Yet direct push is widely feared in production and remains poorly understood at microarchitectural depth. Existing studies either evaluate low-concurrency single-core PIO~\cite{liew2025hotos,kalia2016rdma} or diagnose inbound NIC-to-host DMA congestion in the host interconnect~\cite{agarwal2023hostcc,vuppalapati2024hostnet,liu2023hostping}. What happens when multiple host cores sustain line-rate posted MMIO writes---or when the target device momentarily slows down, drops descriptors, or polls the incoming ring aggressively?

Through systematic hardware-counter measurements across two distinct real-silicon targets sharing the same Sapphire Rapids socket---an \emph{off-chip} PCIe Gen5\,$\times$4 SmartNIC (CN106 \texttt{BAR4}, 32,768 cache lines on \texttt{IIO1}) and an \emph{on-die} accelerator (Intel DSA \texttt{wq2.0}, 4\,KiB / 64-line portal on \texttt{IIO0})---we show that unpaced MMIO direct push breaks down in two ways that appear separate at the application layer but are in fact opposite ends of the same open-loop backpressure chain (\autoref{fig:arch}):

\begin{enumerate}
  \item \textbf{It silently starves co-located host workloads across a universal \texttt{CHA TOR} cliff.} Posted MMIO writes bypass the host LLC and DDR5 channels, but they still allocate entries in the socket's shared 32-slice Caching and Home Agent Table of Requests (\texttt{CHA TOR}) as \texttt{WCiLF} transactions until accepted by \texttt{M2PCIe}/\texttt{IRP}. Whenever concurrent writers push past the egress ceiling (175.5\,Mpps on PCIe Gen5\,$\times$4 or 143\,Mpps on the on-die DSA portal), backpressure refills \texttt{CHA TOR}. Crossing a universal 330--430-entry cliff starves co-located memory requests on the same socket: 6-core STREAM bandwidth plunges by 69.5\% on CN106 \texttt{BAR4} (at 8 cores) and by 67.0\%--96.7\% on the Intel DSA 4\,KiB portal (at 16--31 cores / 62 threads), while 16\,MiB L3 \texttt{p99} latency jumps by 20$\times$ on \texttt{BAR4} and 7.7$\times$--116$\times$ on DSA. Crucially, two intuitive safeguards fail: (i)~hardware descriptor dropping on DSA (\texttt{SWERROR}, 94.6\% drop rate) still drives \texttt{IRP0} and \texttt{M2PCIe0} to 100\% backpressure; and (ii)~limiting a portal to a single 4\,KiB page (64 lines) only kept 8 cores safe by coincidence (318.6 entries, just below the 330 knee) and collapses completely once 12--32 cores share the page. Moreover, at the 12-core onset, physical DDR5 IMC bandwidth already falls by 22.1\% while Golden Cove's adaptive \texttt{RFO}\,$\to$\,\texttt{ItoM} write-allocate shift masks the drop to an apparent 4.2\% in logical STREAM throughput.

  \item \textbf{It collapses its own throughput by 5$\times$--10$\times$ and spills back into the host.} On the CN106 receiver, unpaced writes collide with active receiver polling at two microarchitectural levels. At Level~1, under default cacheable \texttt{Normal-WB} pages, any active receiver core polling an unwritten frontier slot pulls the line into private L1D/L2; incoming PCIe writes then stall behind snoop invalidations in the shared inbound bridge (\texttt{PEM0}/\texttt{IOBN}), locking aggregate throughput at 13.1--13.5\,Mmsg/s regardless of core count ($13.4/K$ per sub-ring) or ARM read ordering (\texttt{ldar}, \texttt{ldr}, \texttt{ldnp}, \texttt{WFE}, \texttt{DC CIVAC}). Even classic speculative skipping (\texttt{B-Queue}~\cite{wang2013bqueue}) backfires by planting ``snoop mines'' that drop single-core throughput from 34.2 to 13.5\,Mmsg/s. At Level~2, remapping the ring with standard Linux \texttt{MT\_NORMAL\_NC} (\texttt{0x44})---zero \texttt{MAIR\_EL1} edits---hits CN106's 48\,MiB shared L3 (\texttt{TAD}) directly with zero L1/L2 snoops and zero DDR5 reads, scaling to 160.1--174.1\,Mpps with four card cores and 901\,ns unloaded RTT. Yet over-provisioning card cores ($K\,=\,8$) or speculative multi-slot closed-loop polling (\texttt{u4}) still triggers L3 \texttt{EXLMN} frontier collisions that collapse closed-loop throughput to 16.3--22.8\,Mmsg/s and overflow the device's $\sim$387--391-slot posted buffer right back into the host \texttt{CHA TOR} cliff (374--521 entries), directly degrading co-located Socket-1 STREAM by $-$11.4\% to $-$12.6\% and inflating co-located pointer-chase \texttt{p99} by +47.6\% to +61.3\%. Meanwhile, pre-\texttt{IIO} traversal across the 32-\texttt{CHA} mesh reorders un-fenced multi-line writes by 2.2\%--5.9\% (and 96.2\%--99.1\% on multi-engine DSA without batch fences), requiring carefully amortized ordering rather than per-packet \texttt{SFENCE}.
\end{enumerate}

\para{Why pacing is the missing piece.}
Both failures share a single root cause: posted MMIO direct push has no native end-to-end flow control. Once either side allows in-flight posted writes to exceed downstream drain capacity, backpressure propagates upstream across PCIe or on-die \texttt{IIO} into the host's shared \texttt{CHA TOR}, turning a local I/O overload into a socket-wide memory outage. Pacing---bounding both aggregate injection rate and micro-burst concurrency on the host, while keeping device-side polling off the incoming write frontier---closes the loop without forfeiting direct push's cache-bypass and latency advantages.

\para{Contributions.}
This paper makes three contributions grounded in 100\% non-multiplexed hardware-counter measurements across two real-silicon platforms:
\begin{itemize}
  \item \textbf{Dual-target characterization of the host \texttt{CHA TOR} cliff (\S\ref{sec:host_contention}).} We establish the unified 330--430-entry \texttt{WCiLF} collapse law across both off-chip CN106 \texttt{BAR4} and on-die Intel DSA, expose the \texttt{ItoM}-masked 12-core IMC drop, disprove hardware-drop and 4\,KiB-page immunity, and identify the \texttt{(Physical Core WCB, 64\,B line)} serialization key and socket-wide line budget $\sum_i K_i < 200$.
  \item \textbf{Two-level dissection of device self-collapse and ordering hazards (\S\ref{sec:device_collapse}).} We isolate the Level-1 \texttt{Normal-WB} snoop lock (13.4\,Mmsg/s), demonstrate zero-\texttt{MAIR} \texttt{MT\_NORMAL\_NC} (\texttt{0x44}) breakthrough to 174.1\,Mpps with zero DRAM traffic, uncover the Level-2 \texttt{EXLMN} collision and host \texttt{CHA} spillover, and quantify pre-\texttt{IIO} mesh reordering for CPU \texttt{MOVDIR64B} and multi-engine DSA.
  \item \textbf{Two complementary pacing pillars and full-stack silicon validation (\S\ref{sec:design}--\S\ref{sec:eval}).} For stateless portals, static disjoint line-budgeting ($K\,=\,2$--$4$) sustains 142.7--142.9\,Mpps across 8--32 cores and 32--62 SMT threads with 99.2\%--101.5\% STREAM retention and zero software flow control. For stateful rings, coupled rate + micro-burst + credit pacing (\texttt{b\,=\,128}, \texttt{W\,$\le$\,256}) restores 0\% neighbor loss at 8-core 165\,Mpps on \texttt{BAR4} and pairs with single-binary zero-lookahead \texttt{MT\_NORMAL\_NC} (\texttt{u1\_ldpq}) on dedicated tight-polling rings (4.5$\times$--12$\times$ over \texttt{Normal-WB}), default \texttt{Normal-WB} on staggered mixed forwarding loops (+10\%--15\% via hardware prefetch), or coalesced doorbells, directly validated with co-located Socket-1 neighbors. On dual-host 100\,GbE, paced direct push delivers +55\% throughput on an in-memory KV store and +22.5\% on zero-modification upstream Memcached~1.6.45 (Facebook ETC), scales Memcached 3.14$\times$ across 1--4 host cores (+13.1\%/+18.9\% over \texttt{SDP} on \texttt{get8}/\texttt{ETC}), and cuts small-request PCIe traffic by 59.4\%--62.5\% and net active CPU cycles by 27.2\%--32.2\%.
\end{itemize}
